\medskip

\textbf{PARTIE A}

\medskip

Le système d'alarme d'une entreprise fonctionne de telle sorte que, si un danger se présente, l'alarme s'active avec une probabilité de $0,97$. 
La probabilité qu'un danger se présente est de $0,01$ et la probabilité que l'alarme s'active est de \np{0,01465}.

On note $A$ l'évènement \og  l'alarme s'active \fg{}  et $D$ l'évènement \og  un danger se présente \fg . 

On note $\overline{M}$ l'évènement contraire d'un évènement $M$ et $P(M)$ la probabilité de l'évènement $M$.



\begin{enumerate}
\item On représenter les éléments de la situation que l'on connaît par un arbre pondéré.

\begin{center}
\pstree[treemode=R,nodesepA=0pt,nodesepB=3pt,treesep = 1cm, levelsep=3cm,nrot=:U]{\TR{}}
	{
	\pstree[nodesepA=3pt]{\TR{$D$~}\naput{$0,01$}}
		{
		\TR{$A$}\naput{$0,97$}
		\TR{$\overline{A}$}\nbput{\blue $1-0,97=0,03$}
		}
	\pstree[nodesepA=3pt]{\TR{$\overline{D}$}\nbput{\blue $1 -0,01=0,99$}}
		{
		\TR{$A$}\naput{}
		\TR{$\overline{A}$}\nbput{}
		}
	}
\medskip
\end{center}

\item  
	\begin{enumerate}
		\item La probabilité qu'un danger se présente et que l'alarme s'active est:
		
$P\left (D\cap A\right ) = 0,01\times 0,97 = \np{0,0097}$.		
		
		\item On en déduit que la probabilité qu'un danger se présente sachant que l'alarme s'active est:
$P_{A}(D)=\dfrac{P\left (D\cap A\right )}{P(A)}=\dfrac{\np{0,0097}}{\np{0,01465}}
\approx 0,662$.		
		

	\end{enumerate}
	
\item La probabilité que l'alarme s'active sachant qu'aucun danger ne s'est présenté est:

$P_{\overline{D}}(A) = \dfrac{P\left (A\cap \overline{D}\right )}{P\left (\overline{D}\right )} $.

On sait que $P(A)=\np{0,01465}$.

D'après la formule des probabilités totales:
$P(A)=P\left (D\cap A\right ) + P\left (\overline D \cap A\right )$.

On déduit: $P(A)-P\left (D\cap A\right ) = P\left (\overline D \cap A\right )$, donc
$P\left (\overline D \cap A\right ) = \np{0,01465}-\np{0,0097} = \np{0,00495}$.

Donc $P_{\overline{D}}(A) = \dfrac{P\left (A\cap \overline{D}\right )}{P\left (\overline{D}\right )} = \dfrac{0,00495}{0,99}= 0,005$.

On peut compléter l'arbre:

\begin{center}
\pstree[treemode=R,nodesepA=0pt,nodesepB=3pt,treesep = 1cm, levelsep=3cm,nrot=:U]{\TR{}}
	{
	\pstree[nodesepA=3pt]{\TR{$D$~}\naput{$0,01$}}
		{
		\TR{$A$}\naput{$0,97$}
		\TR{$\overline{A}$}\nbput{\blue $1-0,97=0,03$}
		}
	\pstree[nodesepA=3pt]{\TR{$\overline{D}$}\nbput{\blue $1 -0,01=0,99$}}
		{
		\TR{$A$}\naput{$0,005$}
		\TR{$\overline{A}$}\nbput{\blue $1-0,005=0,995$}
		}
	}
\medskip
\end{center}

\item On considère qu'une alarme ne fonctionne pas normalement lorsqu'un danger se présente et qu'elle ne s'active pas ou bien lorsqu'aucun danger ne se présente et qu'elle s'active.

Cette situation est représentée par les événements $D\cap \overline{A}$ et $\overline{D}\cap A$.

La probabilité que l'alarme ne fonctionne pas normalement est donc:

$P\left (D\cap \overline{A}\right )+ P\left (\overline{D}\cap A\right )
=  0,01\times 0,03 + 0,99\times 0,005 = \np{0,00525} < 0,01$.

\end{enumerate}

\medskip

\textbf{PARTIE B}

\medskip

Une usine fabrique en grande quantité des systèmes d'alarme. On prélève successivement et au hasard $5$ systèmes d'alarme dans la production de l'usine. Ce prélèvement est assimilé à un tirage avec remise.
On note $S$ l'évènement \og  l'alarme ne fonctionne pas normalement \fg{}  et on admet que $P(S) = \np{0,00525}$.

On considère $X$ la variable aléatoire qui donne le nombre de systèmes d'alarme ne fonctionnant pas normalement parmi les $5$ systèmes d'alarme prélevés.





\begin{enumerate}
\item On a une répétition de 5 épreuves  assimilées à un tirage avec remise donc  la variable aléatoire $X$ qui donne le nombre de systèmes d'alarme ne fonctionnant pas normalement parmi les $5$ systèmes d'alarme prélevés suit une loi binomiale de paramètres $n=5$ et $p=\np{0,00525}$.



\item La probabilité que, dans le lot prélevé, un seul système d'alarme ne fonctionne pas normalement est:
$P(X=1) = \ds\binom{5}{1} \times \np{0,00525}^1 \times (1-\np{0,00525})^{5-1}
\approx \np{0,0257}$.

\item La probabilité que, dans le lot prélevé, au moins un système d'alarme ne fonctionne pas normalement est:

$P(X\geqslant 1) = 1-P(X=0) = 1- \ds\binom{5}{0} \times \np{0,00525}^0 \times (1-\np{0,00525})^5 = \np{0,0260}$.

\end{enumerate}

\medskip

\textbf{PARTIE C}

\medskip

Soit $n$ un entier naturel non nul. On prélève successivement et au hasard $n$ systèmes d'alarme. Ce prélèvement est assimilé à un tirage avec remise.

On cherche le plus petit entier $n$ tel que la probabilité d'avoir, dans le lot prélevé, au moins un système d'alarme qui ne fonctionne pas normalement soit supérieure à $0,07$.

Comme dans la partie B, la variable aléatoire $Y$ qui donne le nombre de systèmes d'alarme défectueux suit une loi binomiale de paramètres $n$ et $p= \np{0,00525}$.

On cherche $n$ tel que $P(Y\geqslant 1) >0,07$.

$P(Y\geqslant 1) >0,07 \iff 1-P(Y=0)>0,07 \iff 0,93 > P(Y=0)$

$\begin{gathered}P(Y=0)<0,93 
\iff \ds\binom{n}{0}\times \np{0,00525}^0 \times (1-\np{0,00525})^n <0,93
\iff \np{0,99475}^n < 0,93\\
\phantom{P(Y=0)<0,93}
\iff \ln \left (\np{0,99475}^n\right ) < \ln(0,93)
\iff n \ln \left (\np{0,99475}\right )  < \ln(0,93)\\
\phantom{P(Y=0)<0,93}
\iff n> \dfrac{\ln(0,93)}{\ln \left (\np{0,99475}\right )}\end{gathered}$

Or $\dfrac{\ln(0,93)}{\ln \left (\np{0,99475}\right )} \approx 13,8$, donc il faut prélever au moins 14 systèmes d'alarme pour que la probabilité d'avoir, dans le lot prélevé, au moins un système d'alarme qui ne fonctionne pas normalement soit supérieure à $0,07$.


\begin{list}{\textbullet}{Vérification à la calculatrice}
\item Pour $n=13$, on trouve $P(Y\geqslant 1) \approx \np{0,0661}<0,07$.
\item Pour $n=14$, on trouve $P(Y\geqslant 1) \approx \np{0,0710}>0,07$.
\end{list}

\vspace{0.5cm}
